\documentclass[10pt]{beamer}

\usetheme{metropolis}
\usefonttheme[onlymath]{serif}
\metroset{sectionpage=none, progressbar=frametitle, numbering=fraction}

\usepackage[T1]{fontenc}
\usepackage[utf8]{inputenc}
\usepackage{booktabs}
\usepackage{ragged2e}
\usepackage{array}

% Tübingen-ish red accent
\definecolor{uniRed}{RGB}{166, 28, 46}
\setbeamercolor{frametitle}{bg=uniRed}
\setbeamercolor{progress bar}{fg=uniRed}
\setbeamercolor{normal text}{fg=black!85}
\setbeamerfont{title}{size=\large, series=\bfseries}

\title{Dialect-Aware Automatic Speech Recognition}
\subtitle{Research background and PhD vision}
\author{Satiyabooshan Murugaboopathy}
\date{Doctoral Researcher Position 1 · University of Tübingen}

\setbeamertemplate{frame footer}{\insertshortauthor\quad$\cdot$\quad\insertshorttitle}

\begin{document}

\begin{frame}[plain,noframenumbering]
  \maketitle
\end{frame}

%% ================= Slide 1: Background =================
\begin{frame}{Background}
  \begin{columns}[T]
    \begin{column}{0.58\textwidth}
      \small
      \textbf{ML researcher}: low-resource, uncertainty-aware modeling and
      \emph{systematic benchmarking} meeting real deployment.
      \vspace{0.4em}
      \begin{itemize}
        \item \textbf{Voice-agent startup (co-founder)} — production Whisper ASR
              (German), wav2vec 2.0 / HuBERT, CTC, on-device inference.
        \item \textbf{Linköping IAS} — uncertainty-aware pipelines for sparse,
              heterogeneous data.
        \item \textbf{Fraunhofer IMW/ISI} — NLP over noisy corpora; weak supervision.
        \item \textbf{Porsche \& Volkswagen} — short sequence / multimodal learning.
      \end{itemize}
    \end{column}
    \begin{column}{0.40\textwidth}
      \small
      \textbf{Education}\\
      M.Sc.\ Computer Science, Uni Leipzig — thesis: systematic benchmarking of
      text classification.\\[0.2em]
      B.Sc.\ Computer Science.
      \vspace{0.5em}
      \begin{alertblock}{Why this project}
        Native \textbf{German} speaker; keen on \emph{speech as a clinical biomarker}.
      \end{alertblock}
    \end{column}
  \end{columns}
\end{frame}

%% ============ Slide 2: Research — interpretability & calibration ============
\begin{frame}{Research I — Understanding \emph{why} models fail}
  \textbf{Main tool: intervene in representation / calibration space, without new
  labels or retraining.}
  \vspace{0.5em}
  \begin{block}{SHEAR — spurious shortcuts vs.\ sparse autoencoders}
    Decomposes frozen text embeddings into sparse features and \emph{ablates}
    shortcut features with no group labels — works even when target and spurious
    attribute are \textbf{perfectly correlated}.
  \end{block}
  \begin{block}{LUCID — statistically grounded self-correction}
    Diagnoses an LLM rater's \emph{systematic bias} (doubly-robust DSL/PPI),
    feeds it back as text, and lets the rater re-grade itself — no retraining.
  \end{block}
  \begin{center}
    \small \textbf{Thread:} make cheap models \emph{trustworthy} under scarcity.
  \end{center}
\end{frame}

%% ============ Slide 3: Research — multimodal & causal verification ============
\begin{frame}{Research II — Verifying representations under scarcity}
  \begin{block}{Platonic Representations (poverty mapping)}
    Fuses \textbf{satellite vision + LLM/agent text} to predict wealth across
    $\sim$60k African DHS clusters; releases an open multimodal dataset.
    Fusion beats vision-only ($R^2$ 0.77 vs.\ 0.63 OOS).
  \end{block}
  \begin{block}{Agent Verifiers (counterfactual analysis)}
    Benchmarks whether agentic treatment-effect representations are
    \emph{causally admissible} — guarding against retrieval and encoder leakage
    before they enter causal analysis.
  \end{block}
  \begin{center}
    \small \textbf{Thread:} benchmarking is itself a research contribution.
  \end{center}
\end{frame}

%% ============ Slide 4: PhD vision ============
\begin{frame}{PhD Vision — Interpretability meets data-scarce ASR}
  \textbf{Problem:} SOTA speech recognition breaks on 1950--75 Swabian dialect —
  regional variation, spontaneous speech, almost no annotation.
  \vspace{0.5em}
  \begin{itemize}
    \item Adapt \textbf{self-supervised speech models} (wav2vec~2.0 / HuBERT /
          WavLM) to dialect — but \emph{diagnose} the shift, don't just fine-tune.
    \item \textbf{SHEAR-style} sparse decomposition of speech features to find
          dialect / shortcut directions.
    \item \textbf{LUCID-style} calibrated uncertainty to route unsure segments to
          expert (human-in-the-loop) curation.
  \end{itemize}
  \vspace{0.4em}
  \begin{center}
    \small Deliverables: uncertainty-aware dialect ASR + IPA transcription,
    an open benchmark, and clinical validation on the TREND cohort.
  \end{center}
\end{frame}

\end{document}